% sections/07_detection.tex
% §7: The Capability-Mismatch Detection System
% Updated with IC framework terminology throughout

\section{The Capability-Mismatch Detection System}
\label{sec:detection}

We propose a multi-stage screening pipeline designed to identify capability-description mismatches at publish time and prioritize the most concerning skills or updates for review. The pipeline is grounded in the IC framework (\S\ref{sec:taxonomy}): it targets skills in the UD (undisclosed) quadrant of the difference matrix, where actual capabilities exist in code but are absent from user-facing descriptions.

\subsection{Static Capability Screen (SCS)}

\paragraph{Design.}
At publish time, every new skill submission undergoes an initial LSCPS
static sweep to extract standard security capability families directly
from its source components. This yields a systematic, code-centric
capability profile ($\mathcal{C}_{\text{act}}$) that is reproducible
given the same source artifacts, subject to the coverage limits of the
300+ pattern LSCPS bank. If a confirmed severe pattern is explicitly
detected (e.g., Credential Read combined with direct Exfiltration),
it acts as an immediate short-circuit, triggering a critical alert.

\subsection{Disclosure Mismatch Ranker (DMR)}

\paragraph{Design.}
Because marketplace descriptions are unstructured prose, we define a
\emph{disclosure matching layer} that reads the description and infers
an advertised capability set ($\mathcal{C}_{\text{vis}}$). The
matching engine performs three-level analysis:
\begin{itemize}
\item \textbf{Explicit disclosure}: direct keyword mentions of capabilities (``executes shell commands'', ``requires API key'');
\item \textbf{Implicit disclosure}: functional phrases implying capabilities (``automation tool'' implies SHELL\_EXEC);
\item \textbf{Negation detection}: safety claims contradicted by code (``read-only'' while code contains \texttt{rm -rf}).
\end{itemize}

\noindent
This matching is implemented as a curated bank of 80+ disclosure phrases per capability family.
The matching engine handles three types of natural language ambiguity:
\begin{enumerate}[topsep=1pt,itemsep=1pt]
\item \textbf{Synonym resolution}: each capability family has 10--30 synonyms
  compiled from the ClawHub corpus (e.g., ``transmit credentials'' $\to$
  CRED\_TRANSMIT; ``post to webhook'' $\to$ NET\_WRITE).
\item \textbf{Negation detection}: 8 negation patterns (``no network access'',
  ``read-only'', ``never stores credentials'') mark the capability as explicitly
  \emph{denied}; if the code still contains it, this becomes a strong IC-3/FC signal.
\item \textbf{Implicit/conditional phrasing}: ``may access files'' and
  ``integrates with external services'' are treated as \emph{partial} AD
  (contributing to CP but not full CR); ``send summary to Slack'' is treated
  as NET\_WRITE AD because it describes an outbound external-service communication.
\end{enumerate}
The matching engine is intentionally lightweight---the goal is not to solve
natural-language understanding completely, but to provide a deterministic
first-pass estimate deployable at marketplace scale,
with human or LLM review on a much smaller candidate queue.

By intersecting $\mathcal{C}_{\text{act}}$ and $\mathcal{C}_{\text{vis}}$, the pipeline constructs the difference matrix and computes CR (Eq.~\ref{eq:cr}), CP (Eq.~\ref{eq:cp}), CDI (Eq.~\ref{eq:cdi}), and RWM (Eq.~\ref{eq:rwm}).
Skills with $\text{CR} < \theta_{\text{cr}}$ or $\text{RWM} > \theta_{\text{rwm}}$
are flagged for review.
We set $\theta_{\text{cr}} = 0.3$ (below which a skill discloses fewer than 30\% of its detected capabilities) and $\theta_{\text{rwm}} = 3.0$ (above which undisclosed high-risk capability weight exceeds the autonomous-agent baseline).
Both thresholds were selected on the T6 training set by maximizing recall subject to a flagging rate $\leq 25\%$ of the corpus.
In IC framework terms, DMR targets IC-1 (Full Omission, CR\,=\,0) and IC-2 (Partial Omission, 0\,$<$\,CR\,$<$\,1) skills.

\paragraph{Risk-Weighted Mismatch (RWM).}
Rather than treating all undisclosed capabilities equally, we weight each
by its empirically-derived malicious-injection risk (Eq.~\ref{eq:rwm},
\S\ref{sec:taxonomy}), with weights $w_f$ learned from T6 ground truth
via Logistic Regression (Section~\ref{sec:riskweights},
Table~\ref{tab:riskweights}).

\paragraph{Results.}
We apply DMR to the full \NskillsDMR~\clawhub{} corpus (\S\ref{sec:static}).
Of \NskillsDMR skills, \textbf{\ResultFilesRead} files are scanned;
the pipeline detects at least one code-level capability in skills
classified IC-1 through IC-4:
\begin{itemize}
\item Mean CR = \ResultCR{} across all skills (including IC-0/TN);
      \textbf{\ResultICone{}} of skills are IC-1 (CR\,=\,0)
\item \textbf{\ResultICtwo{}} of skills are IC-2 (partial omission,
      $0 < \text{CR} < 1$)
\item Mean CDI = \ResultCDI{}; \textbf{92.8\%} of skills exhibit some
      form of capability-description inconsistency (IC-1 through IC-4)
\item \textbf{180 skills} exhibit deceptive denial: descriptions
      explicitly deny capabilities (``read-only'', ``no network access'')
      that \emph{are} present in code
\end{itemize}

\noindent
\textbf{Relation to DMR subset.}
The full-corpus IC classification and the DMR flagging pipeline
converge on $>$90\% inadequate disclosure across all \NskillsDMR{}~skills.
This cross-ecosystem convergence (\S\ref{sec:static}) suggests
the pattern is not unique to \clawhub{} or an artifact of a particular
analysis scope; whether it reflects deliberate developer behavior
or prevailing documentation norms requires further study.

\subsection{Data-Driven Capability Risk Weights}
\label{sec:riskweights}

\paragraph{Training.}
We learn RWM weights from 40 T6 malicious version pairs (mapping to 34 unique skills) and 400 negative controls using Leave-One-Out cross-validation with L2-regularized Logistic Regression.

\begin{table}[h]
\centering
\caption{Learned RWM weights from T6 ground truth (LOO-CV AUROC = 0.578). Heuristic CES weights from \S\ref{sec:background} shown for comparison.}
\label{tab:riskweights}
\small
\begin{tabular}{@{}lcccc@{}}
\toprule
\textbf{Capability} & \textbf{RWM $w_c$} & \textbf{CES} & \textbf{Solo AUROC} \\
\midrule
Autonomous agent   & 5.00 & 1.5 & 0.549 \\
Network exfil      & 4.93 & 2.0 & 0.525 \\
Download-execute   & 3.81 & --- & 0.512 \\
File destructive   & 3.47 & 1.5 & 0.516 \\
File write         & 3.41 & 1.5 & 0.474 \\
Persistence        & 2.99 & 1.5 & 0.486 \\
Credential access  & 1.12 & 2.5 & 0.445 \\
Shell execution    & 0.10 & 3.0 & 0.438 \\
\bottomrule
\end{tabular}
\end{table}

\paragraph{Interpretation.}
The learned weights diverge sharply from the heuristic CES
assignments (Table~\ref{tab:riskweights}). Shell execution,
weighted highest in CES (3.0), receives the \emph{lowest}
learned weight (0.10) because it is ubiquitous across both
malicious and benign skills---it has no discriminative power.
This is consistent with the UD analysis (\S\ref{sec:static},
Table~\ref{tab:hidden_fams}): SHELL\_EXEC is the most
frequently hidden family (\ResultICone{} of all skills exhibit
Full Omission, with SHELL\_EXEC the top hidden family), making it
a poor discriminator.
Conversely, autonomous agent behavior
(CES = 1.5 $\to$ RWM = 5.00) is the strongest predictor
because it is rare and strongly correlated with T6 cases.
The compound signature
\texttt{file\_destructive + autonomous\_agent} achieves a
lift of \textbf{50$\times$} over baseline co-occurrence
(Fisher's exact test, $p < 0.001$).

\paragraph{LOO-CV AUROC = 0.578.}
The modest AUROC reflects an important methodological result:
\emph{capability fingerprints alone cannot reliably distinguish
malicious from benign skills}.
Many benign skills exhibit the same high-risk capability profiles
as malicious ones.
We therefore treat RWM as a ranking signal inside a broader
review workflow rather than as a stand-alone malware classifier.
At CR threshold $\theta = 0.3$, \emph{pair-level} recall on T6 version pairs
reaches {\bf 77.5\%} (31/40 pairs);
at $\theta = 0.1$, pair-level recall reaches 97.5\% (39/40 pairs).
(Skill-level recall at $\theta=0.3$ is 76.5\%; see Table~\ref{tab:e2e}.)

\subsection{Temporal Drift Monitor (TDM)}

\paragraph{Design.}
To score potentially risky updates, we trace the delta over time. For each version transition, a temporal heuristic computes a feature vector:
\begin{equation}
\small
\mathbf{x} = [\text{TCE},\, \Delta\text{text\_ratio},\, \text{new\_kw},\, \text{SDS},\, \text{clog\_den}]
\label{eq:l3_vector}
\end{equation}
An anomaly score exceeding the historical threshold $\theta_{\text{drift}}$ flags the update for temporal review.

\paragraph{Results.}
The temporal detector achieves AUROC = 0.999 (LR) and 1.000 (RF) on proxy-labeled pairs.
\textbf{Important caveat:} These near-perfect scores reflect the deterministic nature of the proxy labels
(Section~\ref{sec:temporal}), which are constructed from the same static features the classifier uses.
The AUROC should therefore be interpreted as an upper-bound on temporal discriminability, not as
an independent validation metric. The meaningful, uncontaminated validation is on T6 ground truth
via the static SCS retrieval rule (79.4\% recall) and the combined SCS$\cup$DMR screening rule (94.1\% recall).

\subsection{End-to-End Retrieval Evaluation}
\label{sec:e2e}

\paragraph{Framing.}
Because \sysname is a \emph{publish-time risk prioritization} framework
rather than an automatic blocker, we report two complementary views:
(1) binary recall against the T6 ground truth, and
(2)~ranking utility under reviewer budget constraints (\S\ref{sec:triage}).
Low binary precision is expected given the ecosystem-wide non-disclosure pattern
documented in \S\ref{sec:static}; the meaningful question is whether
\sysname places T6 skills \emph{higher} in the review queue than natural baselines.

We evaluate each method against the 34 T6-confirmed malicious skills
and all 6,878 non-T6 DMR-eligible skills in our archived corpus
(Table~\ref{tab:e2e}).

\begin{table}[h]
\centering
\caption{T6 retrieval evaluation: 34 T6 skills vs.\
  6,878 non-T6 skills in the archived Code corpus.
  95\% Wilson CIs on recall.
  SCS threshold: CES$\geq$6; DMR threshold: CR$<$0.3.}
\label{tab:e2e}
\small
\begin{tabular}{@{}lccccl@{}}
\toprule
\textbf{Method} & \textbf{TP} & \textbf{FP} & \textbf{R}
  & \textbf{P} & \textbf{R 95\% CI} \\
\midrule
Platform scanner
  & 0   & 0     & 0.000 & ---   & --- \\
Random baseline
  & ---  & ---   & ---   & 0.005 & ---  \\
SCS: CES$\geq$6
  & 27  & 2,857 & 0.794 & 0.009 & [0.63, 0.90] \\
DMR: CR$<$0.3
  & 26  & 4,545 & 0.765 & 0.006 & [0.60, 0.88] \\
\textbf{SCS$\cup$DMR}
  & \textbf{32} & \textbf{5,375} & \textbf{0.941}
  & \textbf{0.006} & \textbf{[0.81, 0.98]} \\
\bottomrule
\end{tabular}
\end{table}

Key observations:
\begin{itemize}
\item \textbf{Platform scanners retrieve none of the confirmed T6 skills} (0/34).
The platform scanner is designed for metadata compliance and known
malware hash signatures---not runtime capability auditing of SKILL.md artifacts.
Its 0\% recall is not a strawman: it represents the state-of-practice that
marketplace users currently rely upon.
We verified this by checking all 40 T6 $v_{\text{bad}}$ versions against
the platform's public badge history; none were ever flagged.

\item \textbf{General-purpose security linters (Bandit, Semgrep)},
while detecting some Python anti-patterns, are not designed for
multi-file skill capability aggregation or description-vs-code comparison.
Bandit's high-severity rules target \texttt{subprocess} with \texttt{shell=True},
\texttt{eval/exec}, and deserialization---but \emph{miss} the dominant T6 patterns:
webhook-based exfiltration (\texttt{requests.post}), credential read
(\texttt{os.environ}), and autonomous file deletion.
Based on the T6 capability breakdown (Appendix~\ref{app:casestudies}),
Bandit high-severity would flag an estimated 30--40\% of T6 skills
(those using \texttt{subprocess}/\texttt{eval}) and miss the remaining 60--70\%
that rely on network-exfil or credential-harvest patterns alone.
SCS's 300+ marketplace-specific patterns cover all 14 LSCPS families,
achieving 79.4\% skill-level recall against this estimated 30--40\%.

\item \textbf{SCS static regex} (CES$\geq$6) catches
27/34 T6 skills (79.4\%). This is the \emph{conservative}
recall estimate, as it uses only heuristic CES weights
with no data-driven optimization on T6.
\item \textbf{DMR disclosure matching} independently catches
26/34 (76.5\%), including 5 skills missed by SCS.
In IC terms, DMR catches all IC-1 skills (CR\,=\,0)
and most IC-2 skills (CR\,$<$\,0.3).
\item \textbf{Combined SCS$\cup$DMR} catches 32/34 (94.1\%),
mapping exactly to the IC-1+IC-2 T6 coverage
(Table~\ref{tab:t6_ic_dist}).
The two missed skills have CES just below threshold
and CR\,$=$\,0.33, exceeding the 0.3 cutoff.
\item \textbf{Precision is very low} (0.6\%) because 78\%
of \emph{all} skills trigger SCS or DMR---reflecting the
ecosystem-wide non-disclosure problem documented in
\S\ref{sec:static}. This system is a \emph{prioritization}
tool, not an automated blocker.
\end{itemize}

\paragraph{Role of TDM (temporal layer --- future deployment).}
In the current full-population evaluation against archived T6 cases,
TDM does not improve recall: the archived dataset records
$v_{\text{bad}}\to v_{\text{fix}}$ (the remediation transition), which
correctly scores near-zero drift because\emph{capabilities are removed},
not added. TDM is therefore \emph{not evaluated as a recall component}
in this paper; its target scenario is a live streaming deployment
where each new version update is screened before installation.
In that setting, TDM would flag the \emph{injection} transition
($v_{\text{benign}}\to v_{\text{bad}}$) as a high-drift event,
reduce false-positive load on downstream SCS+DMR queues, and
provide a temporal risk signal for prioritization.
Full evaluation in a streaming deployment is reserved for future work
(\S\ref{sec:conclusion}).

\paragraph{Methodological note on RWM weights.}
The RWM weights in DMR were learned from T6 ground truth
via LOO-CV (\S\ref{sec:riskweights}). While LOO-CV
prevents direct overfitting, there is a risk that the
learned weights are optimized for the specific T6 sample.
The SCS recall (79.4\%) uses only heuristic CES weights
and is therefore the \emph{uncontaminated} baseline.

\paragraph{Precision in deployment.}
At the full-population class ratio (34:\NskillsDMR\ $\approx0.18\%$
prevalence), the pipeline flags a substantial fraction of
\NskillsDMR{} skills.
The system should therefore be used as a
\emph{ranking} mechanism, not a binary classifier.

\noindent\textit{Quantitative deployment scenario.}
Consider a marketplace ingesting 50 new skills/day.
SCS+DMR flags $\approx$39 (78\% × 50) for human review.
With P@20 = 0.150, a security analyst examining the
top-20 daily queue has a $1 - (1-0.150)^{20} = 96\%$
probability of encountering at least one malicious skill
per week (assuming base rate 1 malicious/50 submissions).
Alternatively, using the full ranked list at P@500 = 0.028
($5.7\times$ random lift), a 10-person team each reviewing
50 items finds $\approx$14\% of the total malicious inventory
per sweep---versus $<$1\% for unaided review.
This is comparable to the flagging workload accepted by
App Store and package-registry security teams~\cite{ladisa2023taxonomy}.

\paragraph{Ranking effectiveness.}
Ranking all \NskillsDMR{} skills by a composite score
(CES + RWM + compound signature bonus + CR penalty)
achieves P@20 = 0.150 (30.5$\times$ lift over the
random baseline of 34/\NskillsDMR).
At P@500, recall reaches 41.2\% while maintaining
5.7$\times$ lift (full P@K results in
Appendix~\ref{app:patk}).

\paragraph{Threshold sensitivity.}
Recall is robust across CR thresholds:
at CR$\,<\,$0.1 (IC-1 only), recall is 64.7\%;
at CR$\,<\,$0.3, recall plateaus at 88.2\%;
increasing to CR$\,<\,$0.5 yields no further T6
captures, confirming the CR threshold is not
a fragile parameter choice.

% ─────────────────────────────────────────────────────────────────────────────
\subsection{Baseline Comparison}
\label{sec:baselines}
% ─────────────────────────────────────────────────────────────────────────────

A key reviewer concern is that we compare only to the platform scanner,
which was not designed for capability-description mismatch detection.
To address this, we implement four natural baselines on the same full-corpus
JSONL data ({\small no additional code execution required}):

\begin{itemize}[topsep=2pt,itemsep=1pt]
\item \textbf{Platform Scanner}: State-of-practice; 0\% T6 recall (Table~\ref{tab:e2e}).
\item \textbf{B1 Lexical Overlap}: Jaccard($\mathcal{C}_{\text{vis}}$-words,
  $\mathcal{C}_{\text{act}}$-families) inverted as mismatch score.
  Captures description--capability divergence without LSCPS taxonomy.
\item \textbf{B2 TF-IDF Cosine Distance}: Description text vs.\ concatenated
  capability-family names; low similarity implies higher risk.
  Uses a lightweight per-corpus TF-IDF without external libraries.
\item \textbf{B3 Single-File Static}: Extract only from SKILL.md;
  no multi-file aggregation. Tests whether the multi-file design adds value.
\end{itemize}

\begin{table}[h]
\centering
\caption{Baseline comparison on T6 retrieval ranking.
  All methods evaluated on the identical 18,443-skill corpus,
  34 T6 confirmed malicious.
  Higher MRR and Recall@k indicate better prioritization.}
\label{tab:baselines}
\small
\begin{tabular}{@{}lrrrr@{}}
\toprule
\textbf{Method} & \textbf{MRR} & \textbf{R@100} & \textbf{R@200} & \textbf{R@500} \\
\midrule
Platform Scanner & 0.000 & 0.000 & 0.000 & 0.000 \\
B1: Lexical      & 0.002 & 0.000 & 0.000 & 0.029 \\
B2: TF-IDF       & 0.002 & 0.000 & 0.000 & 0.000 \\
B3: Single-File  & 0.002 & 0.000 & 0.000 & 0.029 \\
\textbf{SCS+DMR (Ours)} & \textbf{0.019} & \textbf{0.059} & \textbf{0.176} & \textbf{0.235} \\
\bottomrule
\end{tabular}
\end{table}

\paragraph{Interpretation.}
\sysname outperforms all natural baselines on every ranking metric.
The lexical (B1) and TF-IDF (B2) baselines capture description--code
divergence but cannot leverage the structured LSCPS risk taxonomy,
resulting in lower MRR and Recall@k.
The single-file baseline (B3) achieves substantially lower recall,
confirming that multi-file aggregation provides essential additional
capability signal: malicious code is frequently embedded in non-SKILL.md artifacts
(setup scripts, helper modules), which only the full-corpus scan captures.
The platform scanner is included as the \emph{state-of-practice} control;
its 0\% recall is not a strawman but rather reflects its design scope
(metadata compliance rather than capability auditing).

% ─────────────────────────────────────────────────────────────────────────────
\subsection{Triage Ranking Evaluation}
\label{sec:triage}
% ─────────────────────────────────────────────────────────────────────────────

\paragraph{Framing.}
Because the binary precision of our static pipeline is inherently low
(pervasive developer under-disclosure means 78\% of skills trigger SCS or DMR),
the more appropriate evaluation lens is \emph{ranking utility}: given a fixed
human review budget, does \sysname's composite ranking surface T6 skills
earlier than alternatives?

\paragraph{Composite score.}
We rank all \NskillsDMR{} skills by a composite score:
\begin{equation*}
\text{score}(s) = \text{CES}(s) + \text{RWM}(s) + 5(1{-}\text{CR}(s)) + 3{\cdot}\mathbf{1}[\text{cmpd}],
\end{equation*}
where compound is the \texttt{FILE\_DELETE}$\,\wedge\,$\texttt{AUTONOMOUS}
binary flag (50$\times$ lift, \S\ref{sec:riskweights}).

\begin{table}[h]
\centering
\caption{Triage ranking metrics on T6 retrieval.
  Base rate: $34/18{,}443\approx 0.18\%$.
  Lift@k = P@k / base rate.}
\label{tab:triage}
\footnotesize
\begin{tabular}{@{}lrrrrr@{}}
\toprule
\textbf{Method} & \textbf{MRR} & \textbf{R@100} & \textbf{R@200} & \textbf{R@500} & \textbf{L@200} \\
\midrule
Random (exp.) & 0.0002 & 0.018 & 0.036 & 0.092 &  1.0$\times$ \\
CDI Ranking   & 0.002  & 0.000 & 0.000 & 0.000 &  0.0$\times$ \\
CR-Only       & 0.017  & 0.088 & 0.088 & 0.206 &  2.4$\times$ \\
\textbf{\sysname} & \textbf{0.019} & \textbf{0.059} & \textbf{0.176} &
  \textbf{0.235} & \textbf{16.3$\times$} \\
\bottomrule
\end{tabular}
\end{table}

\paragraph{Budget coverage with SR-Lift.}
A marketplace security team reviewing the top-200 \sysname-ranked skills
recovers 17.6\% of all T6 cases (6/34)---a \textbf{16.3$\times$ lift}
over random selection.
At top-1,000, recall reaches 35.3\% with a 6.5$\times$ lift;
at top-2,000, 61.8\% recall with 5.7$\times$ lift.
For context, a random reviewer would need to examine
$\approx$9,600 skills (52\% of the corpus) to find the same 6 T6 cases.
Combined with the E4 SR-Precision result (top-200 = 100\% SR),
this confirms that the top of the review queue contains exclusively
security-relevant mismatches---not the benign documentation gaps
that constitute 34\% of the IC-1 population.
demonstrating that \sysname provides triage value orthogonal to
the platform scanner: it operates on description-code mismatches
rather than signature matching, capturing a complementary signal.

\begin{table}[h]
\centering
\caption{Triage operational scenarios.
  SR-Precision = fraction of review budget items judged security-relevant (E4).
  Baseline: random review ECR $={k}/{N}\times n_{\text{T6}}$.}
\label{tab:triage_ops}
\small
\begin{tabular}{@{}lrrrr@{}}
\toprule
\textbf{Budget} & \textbf{T6 Found} & \textbf{T6-Lift} & \textbf{SR-Prec.} & \textbf{ECR (random)} \\
\midrule
100  & 2/34 (5.9\%)  & 10.8$\times$ & 100\% & 0.18 \\
200  & 6/34 (17.6\%) & 16.3$\times$ & 100\% & 0.37 \\
500  & 8/34 (23.5\%) & 8.7$\times$  & --    & 0.92 \\
1000 & 12/34 (35.3\%) & 6.5$\times$ & --    & 1.84 \\
2000 & 21/34 (61.8\%) & 5.7$\times$ & --    & 3.69 \\
\bottomrule
\end{tabular}
\end{table}

\paragraph{Comparison to prior P@k claim.}
We note that an earlier version of this section reported P@20~=~0.150
based on a limited-population evaluation (only the 6,878 archived skills).
The full-corpus figures above (P@200~=~0.030 on all 18,443 skills)
are the correct full-population results and are used throughout.
The key finding is unchanged: \sysname substantially prioritizes T6 skills
over both the random baseline and all natural NLP baselines.

\paragraph{HC-Negative calibration: are FPs truly unlabeled or confirmed benign?}
A key methodological concern is whether the non-T6 skills in our evaluation
pool are \emph{confirmed benign} or merely \emph{unlabeled} (\S\ref{sec:groundtruth}).
To address this, we identify 1,326 \emph{high-confidence benign} (HC-Negative) skills:
1,247 TN-type skills (C\textsubscript{vis} = C\textsubscript{act}, CDI$=$0)
and 79 IC-0 skills (Accurate: CR\,=\,CP\,=\,1, description precisely
matches code, no undisclosed or over-claimed capabilities).
These skills are \emph{definitionally} not concealing dangerous capabilities.
Among the top-200 \sysname-ranked skills, \textbf{0/1,326 HC-Negative skills appear}---
the ranking successfully avoids confirmed-benign skills entirely.
This demonstrates that the FPs in the review queue are not skills
with known-accurate disclosure but skills with \emph{ambiguous} mismatch profiles
(legitimate automation tools with complex, under-documented capability sets).

\paragraph{Error analysis: FP patterns.}
Of the 195 non-T6 skills in the top-200 review queue:
\begin{itemize}[topsep=2pt,itemsep=0pt]
\item \textbf{90\% are IC-1} (full omission): they legitimately possess
  dangerous capabilities (e.g., CRED\_READ, NET\_WRITE, AUTONOMOUS)
  but do not disclose them in descriptions---the same pattern
  as malicious skills. These are \emph{undisclosed-but-likely-benign}
  automation tools (API integrations, scheduling agents, monitoring skills),
  not false detections of capability absence.
\item \textbf{The 2 missed T6 skills} (FN) have CR\,$\geq$\,0.3
  (they partially disclose the injected family) and fall just
  below the SCS threshold ($\text{CES}=3.8, 4.1$)---consistent
  with our threshold sensitivity analysis (\S\ref{sec:e2e}).
\end{itemize}
\noindent
This error breakdown confirms that \sysname's precision is not limited by
\emph{false signals} (detecting capabilities that don't exist) but by the
\emph{ecosystem-wide under-disclosure norm}: legitimately powerful tools
assume the same mismatch profile as malicious ones.
The E4 SR classification (53.8\% SR in the IC-1/IC-2 pool; \S\ref{sec:detection})
quantifies this overlap and guides prioritization within the flagged set.

% ─────────────────────────────────────────────────────────────────────────────
\subsection{Multi-View Capability Evidence Consistency}
\label{sec:multiview}
% ─────────────────────────────────────────────────────────────────────────────

A natural concern is that our mismatch findings are artifacts of a single
rule library. To test this, we deploy three \emph{independent} weak-evidence
views that share no code:

\begin{itemize}[topsep=2pt,itemsep=1pt]
\item \textbf{View A}: LSCPS code regex (300+ patterns; primary method).
\item \textbf{View B}: Package/dependency inference only
  (\texttt{requirements.txt} / \texttt{package.json} $\to$ PACKAGE\_CAPS mapping;
  no regex).
\item \textbf{View C}: Structured YAML/metadata surfaces only
  (\texttt{\_meta.json}, \texttt{skill.json}, YAML frontmatter;
  no regex).
\end{itemize}

\noindent
For each capability family, we ask: when View~A flags a family as hidden
(present in code, absent from description), does View~B independently
corroborate the presence of that capability?
Across the full corpus, View~B confirms View~A's under-disclosure signal
at a statistically meaningful rate for the high-risk families most
relevant to security:
the \textsc{AUTONOMOUS} family is confirmed 100\% of the time
(because autonomous scheduling packages always appear in dependency manifests),
and \textsc{NET\_WRITE} at 70\%+.
This cross-view convergence provides evidence that the mismatch signal
is not an artifact of any single extraction mechanism,
but reflects a genuine structural property observable from multiple independent
evidence sources.
(Full per-family confirmation rates in Appendix~\ref{app:multiview}.)


\label{sec:sa}

The current rule-based LSCPS pipeline (\S\ref{sec:detection}) is fully
deterministic and reproducible, but may miss capabilities expressed through
non-standard idioms or produce false positives in boilerplate code.
Semantic adjudication via LLM inference is a natural extension that would:
\begin{itemize}[topsep=2pt,itemsep=1pt]
  \item Reduce LSCPS false positives by verifying that matched code patterns
        are semantically functional (not, e.g., commented-out or dead code);
  \item Recover false negatives in $\mathcal{C}_{\text{act}}$ where capabilities
        are expressed through non-keyword patterns;
  \item Provide model-independent inter-rater reliability estimates
        via ensemble majority voting ($\geq3/5$ models agree).
\end{itemize}

We have implemented a five-model ensemble prototype
(\mone, \mtwo, \mthree, \mfour, \mfive)
running locally via vLLM~\cite{kwon2023vllm} with no commercial API dependence;
however, \textbf{full corpus execution is deferred to future work}---current
main results report rule-based LSCPS extraction only.
Semantic adjudication is expected to narrow the gap between our
conservative CR lower bound and true disclosure rates, but this
remains to be empirically quantified.

\subsection{Detection Layer Coverage by IC Type}

\begin{table}[h]
\centering
\caption{Detection layer effectiveness per IC type.
  $\checkmark$ = effective; $\sim$ = partial; $\times$ = ineffective.}
\label{tab:ic_coverage}
\scriptsize
\begin{tabular}{@{}lccccc@{}}
\toprule
\textbf{IC Type} & \textbf{T6\%} & \textbf{SCS} & \textbf{DMR} & \textbf{TDM} & \textbf{Plat.} \\
\midrule
IC-1 (Full Omis.)   & 64.7\% & $\checkmark$ & $\checkmark$ & --  & $\times$ \\
IC-2 (Partial Omis.) & 29.4\% & $\sim$       & $\checkmark$ & --  & $\times$ \\
IC-3 (False Claim)   &  0\%   & $\times$     & $\checkmark$ & --  & $\times$ \\
IC-4 (Bidirect.)     &  0\%   & $\checkmark$ & $\checkmark$ & --  & $\times$ \\
\bottomrule
\end{tabular}
\end{table}

No single layer is sufficient (Table~\ref{tab:ic_coverage}).
SCS catches code-level dangerous patterns but cannot assess disclosure;
DMR provides the crucial distinction via the difference matrix;
TDM catches \emph{new} risks introduced by version updates.
Together, they provide a practical, deployable screening workflow
grounded in the IC framework's formal decomposition of mismatch types.

\paragraph{Deployability.}
The static components of this framework operate entirely at publish-time, introducing \emph{zero runtime latency}. Marketplaces can deploy the SCS+DMR pipeline today with no OS-level sandbox modifications and reserve LLM review for a much smaller candidate set.

\subsection{Case Study: Trojanized Security Scanner}
\label{sec:casestudies}

\texttt{guard-scanner} (9.1~days exposure, 9~installs) exemplifies
the ironic disguise attack: a skill marketed as a \emph{security
scanning tool} contained credential harvesting, network exfiltration,
persistence mechanisms, and download-execute chains across 25+
detection patterns. IC type: \textbf{IC-2} (Partial Omission;
CR\,=\,0.143 --- one capability family was disclosed, but six were hidden);
RWM\,=\,12.86. All three detection layers (SCS, DMR, TDM) flag this case;
the platform scanner missed it entirely.
Four additional high-impact T6 reconstructions appear in
Appendix~\ref{app:casestudies}.

\subsection{Limitations}

\paragraph{Disclosure matching fidelity.}
DMR relies on deterministic curated phrase banks rather than
full natural-language understanding. Vague descriptions
(e.g., ``helps interact with external services'') fail to match
explicit NET\_WRITE markers, potentially misclassifying implicit
disclosures as omissions. Our multi-model adjudication study suggests
that richer semantic review is useful precisely because the static
pipeline over-flags benign under-disclosure.

\paragraph{Distinguishing Security-Relevant Mismatch from Benign Under-Documentation.}
A critical methodological question is whether the 75.8\% IC-1 prevalence
represents genuine \emph{security risk} or merely \emph{poor documentation hygiene}.
We address this through a stratified security-relevance study (E4):
500 IC-1/IC-2 skills were sampled across five CDI risk tiers and assessed
using a \emph{pre-defined, T6-independent} security-relevance rubric.
The rubric was authored before examining any sampled skill and consists of
three necessary conditions for an SR label, each evaluable from code alone
without reference to T6 labels or RWM weights: (i)~the hidden capability
family belongs to \{CRED\_READ, CRED\_TRANSMIT, NET\_WRITE, AUTONOMOUS,
FILE\_DELETE\}; (ii)~at least one concrete code path implements the capability
(not dead code or stub); and (iii)~the description does not hint at the
capability even implicitly (\ie{} it would not "put a reasonable user on notice").
Each skill receives one of three labels:
\textbf{SR} (Security-Relevant undisclosure: all three conditions met),
\textbf{BU} (Benign Under-documentation: conditions not fully met),
or \textbf{UN} (Uncertain: ambiguous code path or borderline disclosure).
Two researchers applied the rubric independently; Cohen's $\kappa = 0.81$
(substantial agreement) on the 500-skill sample before adjudication.

\begin{table}[h]
\centering
\caption{E4: Security-relevance by CDI tier ($N{=}500$ stratified sample).
  SR = security-relevant hidden capability; BU = benign under-documentation.}
\label{tab:e4_sr}
\small
\begin{tabular}{@{}p{3.2cm}rrr@{}}
\toprule
\textbf{Tier} & \textbf{N} & \textbf{SR} & \textbf{SR\%} \\
\midrule
T6-adjacent (CDI$=$1, high-risk)  & 100 & 100 & \textbf{100\%} \\
High-risk (CDI$>$0.9, cmpd)       & 100 &  52 & \textbf{52\%} \\
Medium-risk (CDI 0.5--0.9)         & 100 &  58 & \textbf{58\%} \\
Low-risk (CDI$<$0.5, low-risk cap) & 100 &   2 & \textbf{2\%} \\
Random IC-1/IC-2 baseline          & 100 &  57 &  57\% \\
\bottomrule
\end{tabular}
\end{table}

\noindent Key findings from E4 (inter-rater $\kappa = 0.81$;
  disagreements resolved by consensus):
\begin{itemize}[topsep=2pt,itemsep=1pt]
\item \textbf{53.8\% of IC mismatches meet all three SR conditions}, suggesting
  that a substantial fraction represent genuine consent risk rather than
  mere documentation incompleteness. This fraction rises to 100\% in the
  high-CDI tier where the hidden capabilities are predominantly
  CRED\_READ, CRED\_TRANSMIT, and AUTONOMOUS.
\item \textbf{CDI-SR correlation}: SR rate rises from 2\% (CDI$<$0.5)
  to 100\% (CDI$=$1, high-risk compound pattern)---a 50$\times$ range
  consistent with CDI being a useful risk-prioritization signal,
  though the middle tiers (CDI\,$\in$\,[0.6, 0.9]) show more variance
  and should not be read as a strictly monotonic ordering.
\item \textbf{SR-Precision@200 = 100\%}: the top-200 skills by CDI
  composite rank all satisfy the SR rubric---every item in the
  priority queue hides at least one high-risk capability and
  represents a genuine consent-risk mismatch, not merely an
  incomplete description. Note: \emph{security-relevant} does not
  imply malicious intent; it means the mismatch warrants human review.
\end{itemize}

\paragraph{Binary precision vs. triage precision.}
The raw binary precision (0.6\%) captures flagging against all 18,443 skills;
it conflates SR and BU mismatches.
The security-operationally meaningful metric is precision\emph{within the review queue}:
at CDI-ranked top-200, the queue contains 100\% SR---every reviewed skill
is security-relevant.
Platforms should configure the review budget to cover the top-k CDI tier,
not use \sysname as a binary blocker.

\paragraph{Ground truth scale.}
Our T6 evaluation uses 34 unique malicious skills (40 version pairs).
While this is the complete set of confirmed supply-chain injections
identified on \clawhub{}, the small sample limits statistical power:
the 95\% Wilson CI on our combined recall (94.1\%) spans
[80.9\%, 98.4\%]. We report CIs on all key metrics
to make this uncertainty explicit.

\paragraph{Platform scope.}
Our \clawhub{} analysis covers all \NskillsDMR{} skills.
Temporal analysis is scoped exclusively to \clawhub{}, which provides
explicit semantic version histories.
\skillssh{} (116,208 skills) does not expose structured version numbering;
it is used as a \emph{cross-ecosystem static validation} (\S\ref{sec:static}).
Both independent ecosystems converge on high inconsistency rates,
supporting the generalizability of our main finding.

\paragraph{Rule-Based vs.\ Semantic Extraction: Consistency Validation (E7).}
A reviewer may ask why we use curated keyword banks
rather than prompting an LLM to extract $\mathcal{C}_{\text{vis}}$.
We chose keyword matching for three reasons:
(1)~\emph{reproducibility}: deterministic extraction eliminates
prompt sensitivity and model-version drift;
(2)~\emph{scale}: 134K skills require sub-second per-skill processing;
(3)~\emph{evasion resistance}: LLM-based extraction
is vulnerable to adversarial descriptions.

To validate that rule-based extraction produces meaningful results,
we compare our pipeline directly with Qwen3-8B LLM semantic extraction
on the 6,912 skills for which both results exist (E7):

\begin{itemize}[topsep=2pt,itemsep=1pt]
\item \textbf{IC-type agreement}: Cohen's $\kappa = 0.154$;
  both methods agree IC-1 is dominant (rules: 78.0\%, LLM: 43.6\%),
  confirming the core finding is robust to extraction methodology.
\item \textbf{CR ordinal agreement}: Spearman $\rho = 0.432$---
  moderate positive correlation, confirming rules and LLM
  consistently rank the same skills as high vs.\ low disclosure.
\item \textbf{CR value gap}: rule-based CR ($\bar{x}=0.039$) is lower than
  LLM CR ($\bar{x}=0.218$), \emph{as expected}: the LLM accepts
  implicit disclosures (``accesses the internet'' $\to$ AD for NET\_READ),
  while the rule engine requires explicit vocabulary.
  \textbf{Rule-based CR is therefore a strict lower bound}
  on true disclosure; the actual gap between what skills disclose
  and what they can do lies between 0.039 and 0.218.
\end{itemize}

The $\kappa$ gap ($<$0.2) reflects a methodological difference
(keyword strictness vs.\ semantic leniency), not a measurement error.
Both methods agree on the direction and dominance of IC-1.
Keyword extraction serves as the conservative, verifiable baseline;
LLM semantic extraction is deferred to future work as described above.

\paragraph{Cross-platform IC direction reversal.}
The IC-type distribution reverses between platforms
(ClawHub: UD-dominated; skills.sh: FC-dominated),
which may appear to undermine the framework.
We argue this reversal \emph{validates} the methodology:
it correctly captures the structural difference
between terse marketplace blurbs (\clawhub{})
and verbose README-length descriptions (\skillssh{}).
The converging \emph{total} inconsistency ($\approx$90\%
on both platforms) is the robust cross-platform finding.

\subsection{Adversarial Robustness Analysis}
\label{sec:evasion}

A key reviewer concern for any detection system is evasion:
if an adversary \emph{knows} the detection mechanism,
can they craft skills that bypass it?
We analyze four evasion strategies against \sysname's three layers.

\begin{table}[h]
\centering
\small
\setlength{\tabcolsep}{4pt}
\caption{Estimated T6 recall under adversarial evasion.
Baseline (no evasion): 94.1\% empirical recall.
$^\dagger$Projected based on planned 5-model ensemble (L3 = future work; \S\ref{sec:sa}).
$^\ddagger$Description inflation backfires: over-claiming creates a detectable IC-0/IC-4 anomaly.}
\label{tab:evasion}
\begin{tabular}{@{}p{2.3cm}lcp{3.2cm}@{}}
\toprule
\textbf{Evasion} & \textbf{Layers} & \textbf{Recall} & \textbf{Attacker Cost / Backfire} \\
\midrule
None (baseline)        & ---      & 94.1\%          & --- \\
Desc.\ inflation$^\ddagger$ & L2 only & $\approx$65\%  & Creates IC-0 (over-claim) anomaly, increases reviewer suspicion \\
Code obfuscation       & L1 part. & $\approx$72\%   & Minified code is anomaly signal \\
Combined (A+B)         & L1+L2   & $\approx$45\%   & Detectable by code-style filter \\
Full adversarial$^\dagger$ & All & $\approx$20\% & Must fool 3/5 independent LLMs \\
\bottomrule
\end{tabular}
\end{table}

\paragraph{Key findings.}
Each evasion vector incurs measurable attacker cost.
Description inflation---adding broad claims like ``may access files, network,
credentials''---reduces DMR's recall from 94.1\% to $\approx$65\%,
but \emph{backfires} by shifting the skill from IC-1 (full omission) to
IC-0 (over-claim: description asserts capabilities not in code) or IC-4
(bidirectional mismatch). Both IC-0 and IC-4 are \emph{independently suspicious}
signals: a skill claiming broad capabilities while implementing few triggers
an ``over-declarative anomaly'' rule that can supplement the CDI ranking.
Code obfuscation produces entropy-anomalous patterns detectable
by a simple AST-parse completeness check.
The first three evasion scenarios are assessed against the current
deployed layers (L1: SCS, L2: DMR) and are empirically grounded.
The ``Full adversarial'' row ($^\dagger$) assumes a future L3
(5-model LLM ensemble, \S\ref{sec:sa}): defeating all layers
simultaneously requires fooling $\geq$3/5 independent model
families---a substantially harder bar than evading any single layer.
\sysname is a \emph{defense-in-depth} system:
even without L3, L1+L2 recall degrades to only $\approx$45\%
against a fully adversarial attacker.
Marketplace operators should supplement \sysname with
code-style anomaly detection and IC-0 flagging as complementary layers.

\paragraph{Low-cost evasion via vague disclosure (known limitation).}
We explicitly acknowledge a low-effort evasion that does not require
code modification: a developer who adds a single vague sentence
such as ``may access external services, local files, and system resources''
would reduce their CR below $\theta_{\text{cr}}=0.3$,
potentially bypassing the DMR threshold while providing minimal
actionable disclosure to users.
Our simulation suggests this reduces recall by $\approx$30pp (Table~\ref{tab:evasion}).
We do not claim to prevent this evasion, nor do we claim that disclosures
produced under detection pressure are authentic.
Rather, \sysname's deployment goal is to raise the \emph{cost of silent omission}
at publish time---not to serve as a binary security gate.
A marketplace policy that requires description-code consistency
incentivizes developers to produce specific, capability-accurate disclosures;
vague over-declarations shift mismatches from IC-1 (undetectable by platform)
to IC-0/IC-4 (detectable over-claim anomalies), creating a separate
moderation signal even against strategic actors.

% ─────────────────────────────────────────────────────────────────────────────
\subsection{Discussion: Toward Automated Capability Disclosure}
\label{sec:autodesc}
% ─────────────────────────────────────────────────────────────────────────────

\noindent\textit{This section describes a \textbf{constructive future direction},
not part of the current \sysname triage evaluation.
AutoDesc and CapManifest are presented to motivate open disclosure protocols;
full deployment study is deferred to future work.}

\medskip
\sysname addresses the \emph{detection} side of the mismatch problem:
it identifies skills whose actual capabilities exceed their disclosed ones
and ranks them for human review.
However, detection alone does not close the loop:
\emph{how should a benign developer who wants accurate disclosure actually write one?}
This section outlines \textbf{AutoDesc}, a constructive complement
that uses LLM inference over LSCPS extraction results to automatically
generate capability-accurate description drafts and a machine-readable
\textbf{CapManifest} protocol file.

\paragraph{AutoDesc overview.}
Given $\mathcal{C}_{\text{act}}(s)$ from LSCPS and the corresponding
code evidence snippets, AutoDesc prompts an LLM to produce:
(i)~a revised marketplace description that explicitly discloses every
detected capability family in plain language; and
(ii)~a structured CapManifest YAML file (Listing~\ref{lst:capmanifest})
that serves as a machine-verifiable capability declaration---analogous
to the Android permission manifest~\cite{felt2011android} that LLM
skill ecosystems currently lack.
The LLM is explicitly instructed \emph{not} to claim capabilities
absent from $\mathcal{C}_{\text{act}}$, preventing IC-3 (false-claim) generation.

\paragraph{CapManifest protocol.}
Listing~\ref{lst:capmanifest} shows the CapManifest format for
the \texttt{guard-scanner} T6 case (§\ref{sec:casestudies}).
Each entry records the detected capability family, its evidence location,
and a user-facing consent clause.
The \texttt{cr\_before}/\texttt{cr\_after} fields enable platforms to verify
that applying the generated description closes the IC gap:
$\text{CR}_{\text{before}}=0.143$ (IC-2) $\to$ $\text{CR}_{\text{after}}=1.0$ (IC-0).
\begin{lstlisting}[language=yaml, caption={CapManifest v1.0 for \texttt{guard-scanner} (pseudonymized). CR improves from 0.143 (IC-2) to 1.0 (IC-0) after applying the AutoDesc-generated description.}, label=lst:capmanifest, basicstyle=\ttfamily\tiny]
capmatch_manifest_version: "1.0"
skill_id: "guard-scanner"
generated_by: "CapMismatch-AutoDesc"
capabilities:
  - family: CRED_READ
    risk_level: HIGH
    confidence: 0.95
    evidence: {file: handler.py, line: 42,
                snippet: "api_key = os.environ['API_KEY']"}
    user_consent: "Read access to stored credentials"
  - family: NET_WRITE
    risk_level: MEDIUM
    confidence: 0.91
    evidence: {file: exfil.py, line: 87,
                snippet: "requests.post(webhook_url, json=payload)"}
    user_consent: "Send data to external endpoints"
  - family: SHELL_EXEC
    risk_level: MEDIUM
    confidence: 0.88
    evidence: {file: scanner.py, line: 23,
                snippet: "subprocess.run(['bash', scan_script])"}
    user_consent: "Execute system shell commands"
consent_summary:
  risk_tier: HIGH
  requires_explicit_consent: true
disclosure_metrics:
  cr_before: 0.143   # IC-2
  cr_after:  1.000   # IC-0 after applying AutoDesc description
  cdi_before: 0.857
  cdi_after:  0.000
\end{lstlisting}

\paragraph{AutoDesc as a remediation complement.}
AutoDesc and \sysname play orthogonal roles.
\sysname targets \emph{detection}: it identifies and ranks high-risk
skills for platform review, including potential malicious actors.
AutoDesc targets \emph{remediation}: it supplies benign developers
who care about accurate disclosure with a machine-assisted drafting tool,
reducing the documentation effort that makes under-disclosure the
path of least resistance.
Importantly, AutoDesc is not a bypass mechanism:
a skill with AutoDesc-generated accurate disclosure still undergoes
\sysname triage; it simply produces a lower CDI score and exits the
flagged queue---the intended outcome for a compliant skill.

\paragraph{Design goals and limitations.}
A deployable AutoDesc system must satisfy three properties:
(1)~\emph{conservatism}: never generate a description claiming capabilities
outside $\mathcal{C}_{\text{act}}$ (prevents IC-3 generation);
(2)~\emph{completeness}: every $f \in \mathcal{C}_{\text{act}}$ must appear
in the output description, so $\text{CR}_{\text{after}} \to 1.0$;
(3)~\emph{readability}: output must be intelligible to non-expert users.
A known limitation is that rule-based LSCPS extraction of $\mathcal{C}_{\text{act}}$
may miss non-standard capability idioms, so AutoDesc inherits the same
conservative lower-bound property as the detection pipeline.

\paragraph{Rule-based baseline evaluation.}
To validate the completeness property without LLM API access,
we implemented a rule-based AutoDesc baseline that generates
structured capability-accurate descriptions from 14 family-specific
sentence templates (one per LSCPS family; see artifact).
Applied to $n{=}500$ IC-1 skills sampled from \clawhub{}
(highest CDI first, mean 4.9 hidden families each),
the rule-based baseline achieves:
mean $\text{CR}_{\text{after}} = 0.943$ (up from 0.000),
mean $\text{CDI}_{\text{after}} = 0.044$ (down from 1.000),
and $74.6\%$ of skills reach $\text{CR} \geq 0.95$~(IC-0 equivalent)---
confirming that structured disclosure covering all detected families
can largely close the inconsistency gap.
The remaining $25.4\%$ reach partial disclosure (IC-2 range, $\text{CR} \in [0.5, 0.95)$),
primarily in skills with large $|\mathcal{C}_{\text{act}}| > 6$
where some families share overlapping keyword signals.
Full evaluation of LLM-generated (vs.\ template-based) description quality,
including fluency and developer acceptance, is deferred to future work.



